% No Cartier-divisor induction occurs in this proof.

\begin{proof}
We prove \cref{thm:p171-revised2-complete-resultant}.  First isolate
exactly what is imported from Ducrot.  Ducrot
\cite[Section~2.3.3]{Ducrot97CubeIntersection} shows that each
\(V_i=f_*C_i\) is finite locally free and has the asserted field base
change.  Ducrot \cite[Section~6.3]{Ducrot97CubeIntersection} constructs the
universal sections on \(X\times_SP\).
\cite[Lemma~6.3.1]{Ducrot97CubeIntersection} identifies their intersection
line with \(\mathscr R\), with the positive exponent \(+\delta_i\), and
\cite[Definition~6.3.4]{Ducrot97CubeIntersection} constructs \(r\).  On
each regular-sequence open \(U_\varphi\),
\cite[formula~(17), Corollary~6.1.2, and Section~6.2]{Ducrot97CubeIntersection}
identifies \(r\) with the finite-flat norm of the last universal section.
\cite[Lemma~6.3.2]{Ducrot97CubeIntersection} says that the complement of
\(U=\bigcup_\varphi U_\varphi\) has fibre-local depth at least two.  The
finite-flat norm there is the determinant norm, whose formation commutes
with arbitrary base change
\cite[Lemma~31.18.6, Tag~0BD2]{stacks-project}.  We use no assertion that
Ducrot's construction \cite[Section~6.3]{Ducrot97CubeIntersection}
commutes with arbitrary base change.

Fix a field-valued point \(\operatorname{Spec}k\to S\).  Pull the stage
section and the opens to \(P_k\), and denote them by \(r_k\) and
\(U_k^0\).  Every codimension-one point \(\eta\in P_k\) belongs to
\(U_k^0\).  Indeed, the complement is the field extension of the
corresponding fibre of \(P\setminus U\).  Ducrot's depth-two assertion
\cite[Lemma~6.3.2]{Ducrot97CubeIntersection},
together with regularity of the product of projective spaces, says that
every component of that closed set has codimension at least two; this
property is preserved by a field extension.  Here \(P_k\) is smooth over
\(k\), hence regular by
\cite[Lemma~33.25.3, Tag~056S]{stacks-project}.

Fix such an \(\eta\), choose \(\varphi\) with
\(\eta\in U_{\varphi,k}^0\), and choose an integral affine open

\[
 \eta\in T=\operatorname{Spec}A\subset U_{\varphi,k}^0.
\]

Let \(\mathfrak p\subset A\) be the height-one prime representing
\(\eta\), let \(R=A_{\mathfrak p}=\mathcal O_{P_k,\eta}\), and put
\(d=\dim P_k\).  The first \(n\) universal sections on
\(X_k\times_kT\), in the order \(\varphi\), form a regular sequence on
the entire scheme and cut out a finite flat scheme \(q_T:W_T\to T\).
Write \(L_T\) for the restriction to \(W_T\) of the last universal
invertible sheaf.  Its section induces a map between finite locally free
\(\mathcal O_T\)-modules of the same rank,

\[
 h_T:q_{T*}\mathcal O_{W_T}\longrightarrow q_{T*}L_T.
\]

It is an isomorphism at the generic point of the integral scheme \(T\),
because the generic universal tuple has empty common zero.  Hence
\(h_T\) is injective.  Since \(q_T\) is finite, the original map
\(\mathcal O_{W_T}\to L_T\) is injective as well.  If
\(\mathcal I_T\subset W_T\) is its zero scheme, exactness of finite
pushforward gives

\begin{equation}
 \mathcal C_T:=\operatorname{coker}(h_T)
 =q_{T*}(L_T|_{\mathcal I_T}).
 \label{eq:p171-revised2-cokernel-sheaf}
\end{equation}

Base change to \(R\).  The finite \(R\)-algebra
\(B=\Gamma(W_T\times_T\operatorname{Spec}R,\mathcal O)\) is finite free.
It is semilocal, so the restriction of \(L_T\) is free.  After choosing a
generator, the last section is an element \(a\in B\), and the norm section
is represented by \(\det(m_a:B\to B)\).  The determinant is nonzero, and
the same torsion-free argument makes \(m_a\) injective.  Smith normal form
over the DVR gives (compare
\cite[\emph{Lengths and determinants}, Tag~02QA]{stacks-project})

\begin{equation}
 \operatorname{ord}_\eta(r_k)
 =\operatorname{ord}_R\det(m_a)
 =\operatorname{length}_R(B/aB)
 =\operatorname{length}_R((\mathcal C_T)_\mathfrak p).
 \label{eq:p171-revised2-dvr-length}
\end{equation}

We now calculate the incidence cycle on the affine neighbourhood \(T\),
before localizing its coefficient.  This order is essential: the morphism
\(\operatorname{Spec}R\to P_k\) is not used as a flat pullback in
\cite[Tag~02RA]{stacks-project}.  The projection
\(p_T:X_k\times_kT\to X_k\) is a flat morphism between schemes locally of
finite type over \(k\), so
\cite[Definition~42.14.1 and Lemma~42.14.4, Tag~02RA]{stacks-project}
apply to \(p_T\).  They identify the starting
top-dimensional coherent-sheaf cycle with
\(p_T^*[X_k]_n\), retaining the generic lengths.

Apply the Stacks Project's effective-Cartier identity for coherent-sheaf
cycles \cite[Lemma~42.29.5, Tag~02TB]{stacks-project}, successively on
\(X_k\times_kT\) to the first \(n\) universal sections.  At every stage
the section is regular on the entire preceding quotient.  The
coefficient-by-coefficient computation in the proof of
\cite[Tag~02TB]{stacks-project} therefore yields the associated-cycle
identity

\begin{equation}
 [\mathcal O_{W_T}]_d
 =c_1(\mathcal L_{\varphi(n-1)})\cdots
  c_1(\mathcal L_{\varphi(0)})\cap p_T^*[X_k]_n.
 \label{eq:p171-revised2-first-n-cuts}
\end{equation}

The injectivity proved above permits one last use of the same identity:

\begin{equation}
 [L_T|_{\mathcal I_T}]_{d-1}
 =c_1(L_T)\cap[\mathcal O_{W_T}]_d.
 \label{eq:p171-revised2-last-cut}
\end{equation}

For an integral top-dimensional component \(Z\subset X_k\), let
\(V_i\otimes\mathcal O_Z\twoheadrightarrow C_i|_Z\) have kernel
\(Q_i|_Z\).  The full incidence over \(Z\) is the product over \(Z\) of
the projective bundles of lines in the \(Q_i|_Z\).  Its restriction over
\(T\) is exactly the common zero scheme used in
\eqref{eq:p171-revised2-first-n-cuts}--\eqref{eq:p171-revised2-last-cut}.
Thus, by linearity in
\([X_k]_n=\sum_Z\ell(\mathcal O_{X_k,\eta_Z})[Z]\), those two identities
say

\begin{equation}
 [L_T|_{\mathcal I_T}]_{d-1}
 =j_{\mathcal I}^*\Gamma_k([X_k]_n),
 \label{eq:p171-revised2-incidence-restriction}
\end{equation}

where \(j:T\hookrightarrow P_k\) and \(j_{\mathcal I}\) is the induced
open immersion on the incidence scheme.  Here
\cite[Tag~02RA]{stacks-project} is used
only for the finite-type flat projection \(p_T\) and for open immersions,
never for the local-spectrum morphism
\(\operatorname{Spec}R\to P_k\).

Proper pushforward of coherent-sheaf cycles
\cite[Lemma~42.12.4, Tag~02R3]{stacks-project} and
\eqref{eq:p171-revised2-cokernel-sheaf} give

\begin{equation}
 [\mathcal C_T]_{d-1}
 =q_{T*}[L_T|_{\mathcal I_T}]_{d-1}
 =j^*q_{k*}\Gamma_k([X_k]_n).
 \label{eq:p171-revised2-open-pushforward}
\end{equation}

The last equality is also elementary: proper pushforward commutes with
restriction to an open Cartesian square, as is checked on each integral
cycle generator from
\cite[Definitions~42.12.1 and~42.14.1]{stacks-project}; the relevant
function-field degree is unchanged after restriction.  Finally, by the
definition of a coherent-sheaf cycle \cite[Tag~02QV]{stacks-project}, the
coefficient
of the height-one point \(\eta\) in the left side of
\eqref{eq:p171-revised2-open-pushforward} is
\(\operatorname{length}_{A_\mathfrak p}((\mathcal C_T)_\mathfrak p)\).
Combining this with \eqref{eq:p171-revised2-dvr-length} identifies the
coefficient at \(\eta\) of the two sides of
\eqref{eq:p171-revised2-divisor-cycle}.  This includes all residue-field
degrees through \cite[Tag~02R3]{stacks-project} and all generic component
lengths through \cite[Tag~02QV]{stacks-project}.

Since every codimension-one point was arbitrary and \(P_k\) is regular, a
section of a line bundle has divisor determined by these DVR orders.  This
proves \eqref{eq:p171-revised2-divisor-cycle}.  The cycle is precisely the
one in Chen--Moriwaki
\cite[Definition~1.6.3]{ChenMoriwaki21AdelicIntersection}; lower-dimensional
and embedded
components do not occur in \([X_k]_n\), while every top-dimensional generic
nilpotent thickness is the displayed local length.

It remains to determine the scalar and orientation.  Pass faithfully flatly
to a purely transcendental extension \(K/k\), hence to an infinite field.
Set \(Y_0=X_K\).  For \(0\leq j<n\), recompute the full finite set
\(\operatorname{Ass}(Y_j)\).  For each associated point, sections in the
complete basepoint-free system of \(C_j\) that vanish there form a proper
linear subspace.  A finite union of proper linear subspaces cannot exhaust
the section space over \(K\); choose \(t_j\) outside their union and put
\(Y_{j+1}=V(t_j)\subset Y_j\).  Multiplication by \(t_j\) is injective on
all of \(\mathcal O_{Y_j}\), so \(Y_{j+1}\) is an effective Cartier
divisor on the entire possibly nonpure and nonreduced scheme \(Y_j\).
Embedded and lower-dimensional components are included in the associated
point avoidance: they are either missed or cut, and none is discarded.
The choice is repeated on the new quotient rather than inferred from an
earlier associated-prime calculation.

After \(n\) steps, \(Y_n\) has dimension at most zero.  Indeed, on every
surviving positive-dimensional projective component the restriction of a
very ample line bundle cannot have a nowhere-vanishing section, so a regular
section cuts that component in dimension.  Choose \(t_n\) nonvanishing at
every point of the finite support of \(Y_n\).  The total intersection is
empty.

Let \(t=(t_0,\ldots,t_n)\).  Ducrot's successive positive restriction maps
and final norm \cite[Sections~6.1--6.2]{Ducrot97CubeIntersection} define
\(e_t\in I_{X_K/K}(C_0,\ldots,C_n)\setminus\{0\}\).  Write
\(r_K=e_tF_t\).  By construction \(F_t(t)=1\).  Let \(R_t\) be the
Chen--Moriwaki resultant normalized by the literal equality \(R_t(t)=1\),
as in \cite[Definition~1.6.7]{ChenMoriwaki21AdelicIntersection}.  The
divisor identity gives
\(\operatorname{div}(F_t)=\operatorname{div}(R_t)\).  Hence
\(F_t/R_t\) is a global unit on the integral projective coefficient product
\(P_K\).  Since \(\Gamma(P_K,\mathcal O_{P_K})=K\), that ratio is a scalar
in \(K^\times\); evaluation at \(t\) makes the scalar literally \(1\).
Thus \(F_t=R_t\), proving
\eqref{eq:p171-revised2-positive-normalization} over \(K\).  Equality
descends to \(k\) because scalar extension of a vector space is injective.

It remains to separate the rank-one, zero-dimensional endpoint.  If
\(n=0\), Ducrot \cite[Proposition~6.1.1(3)]{Ducrot97CubeIntersection}
identifies the intersection line with
the positive finite-flat norm, and
\eqref{eq:p171-revised2-dvr-length} is the determinant of multiplication on
the original finite algebra, as in Chen--Moriwaki
\cite[Example~1.6.4]{ChenMoriwaki21AdelicIntersection}.  If in
addition \(\dim_kV_{0,k}=1\), very ampleness makes the complete morphism a
closed immersion
\(X_k\hookrightarrow\mathbf P(V_{0,k}^\vee)=\mathbf P^0_k\), so it forces
\(X_k=\operatorname{Spec}k\).  Choose a nonzero basis section \(v\).  On
the affine coefficient line the universal section is \(a_0v\): for a
nonzero tuple \(a_0\ne0\), and the determinant norm polynomial is the
single linear coordinate \(a_0\).  On \(\mathbf P^0_k\), this coordinate
is nowhere vanishing as a section of \(\mathcal O(1)\), so the incidence
and projective divisor are empty.  Framing by the value at the unique
projective tuple makes the normalized positive norm the unit \(1\), not a
section of a dual norm line.
\end{proof}

\begin{proof}
We prove \cref{thm:p171-revised2-one-stage-bridge}.  The maps on
symmetric powers are ordinary contravariant substitution of linear forms,
so kernels of the \(a_i\) cause no algebraic problem.  We first prove the
comparison after passage to a field \(k\), without constructing a
projective map between coefficient spaces.

Let \(Z\) be an integral top-dimensional component of \(X_k\).  The
complete resultant polynomial, after affine substitution by the \(a_i\),
vanishes at \(e=(e_0,\ldots,e_n)\) exactly when the evaluated sections
\(a_i(e_i)\) have a common zero on \(Z\).  This is the degeneracy locus for
the closed immersions defined by the \(E_i\).  Chen--Moriwaki
\cite[Proposition~1.6.2 and Definition~1.6.3]{ChenMoriwaki21AdelicIntersection}
say that this locus is an irreducible
hypersurface.  The substituted polynomial is nonzero because \(n+1\)
general members of these basepoint-free ample systems have empty common
intersection.  It remains homogeneous in block \(i\) of the unchanged
positive degree

\[
 \delta_i(Z)=(C_0\cdots\widehat C_i\cdots C_n)_Z.
\]

Equality of zero sets gives a positive power of the irreducible equation;
comparison in any block gives exponent one.  Hence the substituted complete
polynomial is a scalar times the \(E_i\)-resultant of \(Z\).

For the fundamental cycle,
\cref{thm:p171-revised2-complete-resultant} and Chen--Moriwaki
\cite[Definition~1.6.3]{ChenMoriwaki21AdelicIntersection} factor the
complete polynomial as the product over \(Z\),
with exponent \(\ell(\mathcal O_{X_k,\eta_Z})\).  Affine substitution
distributes over this product, so the same generic lengths occur for the
\(E_i\)-resultants.  After a purely transcendental extension if needed,
choose a tuple from the \(E_i\) by successively avoiding all associated
points of every new quotient, as in the preceding proof.  The first \(n\)
members then form a regular sequence on all of \(X_k\), and the total
intersection is empty.  Evaluation of the substituted line-valued tensor
is the same positive norm element \(e_t\), so its scalar polynomial takes
the literal value \(1\).  Chen--Moriwaki's normalization
\cite[Definition~1.6.7]{ChenMoriwaki21AdelicIntersection} fixes the remaining
scalar.  Equality descends from the faithfully flat field extension.

Now let \(A_0\subset A\) be the one chosen stage and let
\(A_0\to A\to k\) be any factorization through a field.  Associativity of
tensor product says that the \(k\)-fibre of the substituted \(A_0\)-tensor
pulled to \(A\) is obtained as follows: first take the \(k\)-fibre of the
complete \(A_0\)-section, and then apply the linear substitutions
\((a_i)_k\).  Apply
\cref{thm:p171-revised2-complete-resultant} directly to the field
map \(\operatorname{Spec}k\to\operatorname{Spec}A_0\).  It identifies the
first object with the complete fundamental-cycle resultant over \(k\).
Apply the affine zero-set and block-degree argument just proved over \(k\)
to identify its substitution with the finite-system Chen--Moriwaki
resultant \cite[Definition~1.6.3]{ChenMoriwaki21AdelicIntersection}.  This
argument does not pull an \(A_0\)-regular open through
\(A\to k\), and it does not use a projective map associated to an
\(a_i\) having a kernel.  It proves fibre correctness without proving model
independence.  The intersection line, unlike Ducrot's construction
\cite[Section~6.3]{Ducrot97CubeIntersection}, is
the alternating finite difference of determinants of perfect complexes
\(Rf_*\bigotimes_{i\in I}C_i\); proper perfect base change
\cite[Tag~0B91]{stacks-project} gives its arbitrary base change.

For completeness, consider \(n=0\) and \(\dim_kV_{0,k}=1\).  The first
proof gives \(X_k=\operatorname{Spec}k\), and the surjective evaluation
makes \((a_0)_k:E_{0,k}\to V_{0,k}\) a nonzero linear map.  If \(z_0\)
is the affine coordinate on \(V_{0,k}\), the complete affine resultant is
\(z_0\), and its substitution is the nonzero linear form
\(z_0\circ(a_0)_k\).  Its projective zero divisor is exactly
\(\mathbf P(\ker((a_0)_k))\), the incidence divisor for the embedded point;
it is empty when \(E_{0,k}\) also has rank one.  Choosing
\(e\in E_{0,k}\) with \((a_0)_k(e)=v\) makes the normalized scalar value
literally \(1\), in agreement with the positive finite norm.  This is
separate from the complete rank-one projective coefficient space, whose
divisor is always empty.
\end{proof}

\begin{proof}
We prove \cref{cor:p171-revised2-veronese-resultants}.  Descend the
finite presentation, the line bundles, sections, evaluation maps, and the
original closed immersions to one finitely generated
\(\mathbf Z\)-subalgebra \(A_0\subset A\), enlarging it until the model is
flat.  The needed limit statements are
\cite[Tags~01ZM, 01ZR, 0B8W, and~05LY]{stacks-project}.  Since \(A_0\subset A\)
and \(A\) is integral, \(A_0\) is a Noetherian domain.  Proper flat finite
presentation and \cite[Tag~0D4R]{stacks-project} give constant fibre
dimension \(n\).

Relative Serre vanishing gives \(p_0\) such that \(B_{i,0}^p\gg0\) for all
\(i\) and \(p\geq p_0\).  The all-monomial system is the composite of the
original closed immersion with the degree-\(p\) Veronese closed immersion.
Apply \cref{thm:p171-revised2-one-stage-bridge}.  Ducrot
\cite[Lemma~6.3.1]{Ducrot97CubeIntersection} computes the block degree as

\[
 (B_{0,k}^p\cdots\widehat{B_{i,k}^p}\cdots B_{n,k}^p)
 =p^n(B_0\cdots\widehat B_i\cdots B_n),
\]

including the generic multiplicities of the fibre.  This proves the
corollary.
\end{proof}
